Finally, we consider a (mostly) theoretical issue.
If we look to our master formula, \cref{eq:fact1}, we recall it is a (functional) product between two unphysical objects: the coefficient function, $\vb C$, and the PDF, $\vb f$.
Since only the product, the cross section, $\sigma$, is a physical object, there is some freedom in the definition of $\vb C$ and $\vb f$.
We discuss this freedom in \cref{sec:sv} for factorization scale variation schemes, but here we take a more fundamental point of view: how do we define PDFs in the first place?

This freedom is referred to as PDF schemes and DIS is playing a key role here.
We find \cref{eq:LOC} at LO accuracy and \cref{eq:NLOCq,eq:NLOCg} at NLO in the $\overline{\mathrm{MS}}$ scheme, which is the most common PDF scheme.
In the DIS scheme~\cite{Diemoz:1987xu} we define for the coefficient functions of the structure function $F_2$
\begin{equation}
    C_q^{\mathrm{DIS}}(z,Q^2) = C_{\bar q}^{\mathrm{DIS}}(z,Q^2) = e_q^2 \delta(1-z)\,\forall q \qq{and} C_g^{\mathrm{DIS}}(z,Q^2) = 0 \label{eq:CDIS}
\end{equation}
which holds to all perturbative orders.
Since the structure function has to remain invariant of course, i.e.\
\begin{equation}
    F_2(x,Q^2) = \left(\vb C^{T,\overline{\mathrm{MS}}}(Q^2) \otimes \vb f^{\overline{\mathrm{MS}}}(Q^2) \right)(x)
        = \left(\vb C^{T,\mathrm{DIS}} \otimes \vb f^{\mathrm{DIS}}(Q^2) \right)(x)\,, \label{eq:F2DIS}
\end{equation}
this implies, first, that also the PDFs are different and, second, as this must be true at any $Q^2$ also the PDF evolutions are different.
In other words: the PDF scheme defines what we call a quark or a gluon.
In practice all modern PDFs are determined in the $\overline{\mathrm{MS}}$ scheme and thus the DIS scheme is not of any practical relevance, but it can provide some practical advice based on theoretical considerations.

Before making explicit use of the DIS scheme, we review the interplay with the previous content.
\begin{itemize}
    \item We need all possible DIS combinations (NC vs.\ CC) to define all quark combinations, see \cref{sec:diverse}.
    \item The prescription only defines quarks, but not the gluon, which is explicitly defined to be not contributing.
        Instead, one needs an additional definition of the gluon distribution, for which various options exist~\cite{Diemoz:1987xu,Soar:2009yh,Altarelli:1998gn,Candido:2023ujx}.
    \item \Cref{eq:F2DIS} explicitly defines $F_2$ and thus any observable other then $F_2$, e.g.\ $F_L$ from \cref{sec:xs}, becomes more complicated, since, e.g.\ in this case the sum has to be invariant.
        Of course this also applies to massive coefficient functions from \cref{sec:hq} and while the general construction, e.g.\ for FONLL, is the same, the ingredient become more complicated.
    \item Splitting functions $\vb P$ depend on the PDF scheme and thus (almost) all complications discussed in this work just appear there, see \cref{sec:resum}.
    \item TMC, see \cref{sec:TMC} are universal kinematic corrections which apply to any leading-twist expression and are thus independent of the PDF scheme.
        Thus, the naive simplicity of \cref{eq:CDIS} is destroyed when considering TMC, as we must for any real-life DIS comparison.
\end{itemize}

Considering PDF schemes other then the default $\overline{\mathrm{MS}}$ scheme can sometimes be advantages~\cite{Delorme:2026vln}, even if just for theoretical arguments as we illustrate next.
The difference between \cref{eq:CpQCD} (implicitly in $\overline{\mathrm{MS}}$ scheme) and \cref{eq:CDIS} is, by construction, a perturbative object and thus the scheme transformation between these two schemes is so as well.
In particular, we write
\begin{equation}
    \vb f^{\mathrm{DIS}}(\xi,Q^2) = \left(\vb S^{\mathrm{DIS}\leftarrow\overline{\mathrm{MS}}}(Q^2) \otimes \vb f^{\overline{\mathrm{MS}}}(Q^2)\right)(\xi)
\end{equation}
with a perturbative operator $\vb S^{\mathrm{DIS}\leftarrow\overline{\mathrm{MS}}}(Q^2)$.
In this specific case we simply have
\begin{equation}
    \vb S^{\mathrm{DIS}\leftarrow\overline{\mathrm{MS}}}(Q^2) = \left(\vb C^{T,\mathrm{DIS}}(Q^2)\right)^{-1} \vb C^{T,\overline{\mathrm{MS}}}(Q^2) \label{eq:DIStrafo}
\end{equation}
where the inversion is to be understood in a functional sense.

Now, since \cref{eq:F2DIS} defines an observable\footnote{or one can generalize it such that it actually does~\cite{Candido:2023ujx}} and \cref{eq:CDIS} only contains positive numbers, this implies that $\vb f^{\mathrm{DIS}}(Q^2)$ must also be positive.
Recall that cross section are physical objects and thus positive by definition, PDFs are not physical objects and their positivity in non-trivial in general.
In the DIS scheme, the two objects (quasi) coincide and thus the positivity of DIS PDFs follows.
However, together with the perturbativity property of $\vb S(Q^2)$ also the positivity of $\vb f^{\overline{\mathrm{MS}}}$ follows as long as $\vb S(Q^2)$ is actually perturbative~\cite{Candido:2023ujx}.
Checking \cref{eq:DIStrafo} explicitly yields the condition that $\overline{\mathrm{MS}}$ PDFs should be positive above $Q^2 \gtrsim \SI{5}{\GeV^2}$, which has practical consequences for PDF fitting.
